\chapter{基于地理信息增强的 GNPR-SID 方法}
\label{cha:method}

本章主要介绍本文所采用的方法框架及其改进思路。本文的研究工作建立在 GNPR-SID 框架基础之上，其核心目标是面向 Next POI 推荐任务，学习适合生成式模型建模的离散语义标识，并利用大语言模型完成下一兴趣点预测。针对原始 GNPR-SID 在地理信息建模方面存在的不足，本文尝试从地理信息注入方式出发，对原始方法进行改进。经过对显式注入与隐式注入两种方案的比较分析，本文最终选择在 Semantic ID Construction 阶段引入地理增强表示，通过 GeoPE 与文本编码器构造包含空间信息的连续 POI 表示，并将其送入 RQ-VAE 学习地理增强的语义 ID，从而提升 POI 表示的空间感知能力。

\section{GNPR-SID 原始框架}
\label{sec:gnprsid_framework}

GNPR-SID 是一种面向生成式 Next POI 推荐任务的方法，其整体框架可以划分为两个核心组件：Semantic ID Construction 模块与 Generative POI Recommendation 模块。前者负责为每个 POI 学习具有语义结构的离散表示，后者则负责利用这些语义 ID 对用户下一步可能访问的兴趣点进行生成式预测。该框架的核心思想是：将传统推荐系统中难以建模的离散 POI ID 替换为具有层次结构和语义相似性的离散 token 序列，使大语言模型能够在统一的 token 空间中学习用户历史行为与目标 POI 之间的映射关系。

在 Semantic ID Construction 模块中，GNPR-SID 首先基于 POI 的多源属性构造连续表示。原始方法将 POI 的语义信息和协同信息进行联合建模，主要整合了四类特征。第一类是基本类别信息，用于描述该 POI 的语义属性；第二类是位置相关信息，原方法将经纬度通过 Google Maps 的 Plus Codes 转换为区域级地理编码，用于提供粗粒度的区域位置特征；第三类是时间信号，通过历史访问时间映射到离散时间槽，以反映 POI 的时间活跃模式；第四类是用户协同信号，包括最频繁访问用户以及共访问模式等，用于从行为协同角度增强 POI 表示。将上述特征分别表示为 $c$、$r$、$t$ 和 $c_u$，则原始 GNPR-SID 将其拼接形成 POI 的连续表示：
\begin{equation}
p_e = \mathrm{concat}(c, r, t, c_u).
\end{equation}

在获得 POI 的连续表示 $p_e$ 后，GNPR-SID 采用 RQ-VAE 对其进行多级残差量化，将连续表示映射为离散语义标识 SID。设模型包含 $L$ 层量化码本，则每个 POI 在各层码本中对应一个 codeword index，最终 SID 可表示为
\begin{equation}
SID = [\mathrm{index}(e^{(1)}), \mathrm{index}(e^{(2)}), \dots, \mathrm{index}(e^{(L)})].
\end{equation}
由于残差量化具有从粗到细逐层编码的特性，语义相似的 POI 往往会共享更长的 SID 前缀。这种层次化离散表示方式为后续生成式模型提供了更适合学习的目标结构。当不同 POI 映射到相同 SID 时，原始方法还会增加额外维度以缓解 collision 问题。

为了学习高质量的离散语义表示，GNPR-SID 通过联合优化重建损失、量化损失和多样性损失对 RQ-VAE 进行训练，其总体目标函数可表示为
\begin{equation}
L_{total} = L_{recon} + \mu L_{quant} + \lambda L_{div},
\end{equation}
其中，$L_{recon}$ 用于约束解码器能够从量化表示中恢复原始输入特征，$L_{quant}$ 用于优化编码器输出与码本向量之间的匹配关系并稳定量化过程，$L_{div}$ 则用于提高码本使用的多样性，避免过多样本集中映射到少量 codeword 上。

在 Generative POI Recommendation 模块中，GNPR-SID 将用户历史访问序列转换为结构化文本提示，并通过 supervised fine-tuning 的方式微调大语言模型，使模型能够根据历史行为序列生成目标 POI 的 SID。设用户 $u$ 的历史访问序列为
\begin{equation}
H_u = \{(p_1,t_1),(p_2,t_2),\dots,(p_n,t_n)\},
\end{equation}
其中 $p_i$ 为第 $i$ 次访问的 POI，$t_i$ 为对应时间戳。原始方法首先将每个 POI 替换为其对应的 SID token，再与访问时间一起组织为输入提示。其输入形式通常可以表示为：
\begin{quote}
\small
\texttt{user\_<uid> visited: <SID> at [time], ..., visited <SID> at [time]. When [time] user\_<uid> is likely to visit:}
\end{quote}
输出目标则为下一兴趣点对应的 SID 序列。通过这种方式，GNPR-SID 将 Next POI 推荐任务重构为一个条件生成问题，并在原论文中基于 LLaMA3-8B 等大语言模型进行微调训练。

从整体上看，GNPR-SID 的优势主要在于利用层次化离散语义表示替代随机 POI ID，使生成式模型能够更自然地学习 POI 之间的语义相似性及用户访问模式。然而，该方法对地理信息的利用仍然偏弱。虽然原始方法在 Semantic ID Construction 阶段考虑了 region-based code 等位置特征，但其空间信息仍然较粗粒度，且未与语义表示形成更深层次的耦合。这也为本文进一步改进提供了空间。

\section{问题分析与地理信息引入思路}
\label{sec:problem_analysis}

POI 推荐任务与一般序列推荐任务的重要区别在于，POI 天然携带显著的空间属性。用户下一步访问的地点并不仅由历史兴趣偏好决定，还受到当前位置、活动半径、区域可达性以及城市空间结构的影响。因此，如果离散语义标识仅由类别、时间和协同信号等特征学习得到，而缺乏对真实地理空间关系的充分建模，那么所得 SID 对 POI 的表达能力仍然是不完整的。

在 GNPR-SID 原始框架中，POI 的空间信息主要通过 region-based code 以较为粗糙的方式加入 POI 连续表示。这种做法能够提供一定区域级别的空间先验，但仍存在两个问题。首先，region code 更接近一种符号化区域标签，其表达能力有限，难以反映 POI 间更加细粒度的地理差异。其次，这种地理信息与 POI 语义表示之间更多是简单拼接关系，而未在表示学习过程中形成更深层次的耦合。因此，最终学得的 SID 往往仍以语义信息为主，对空间分布模式的刻画不够充分。

针对这一问题，本文的核心思路是：通过更合理的方式让大语言模型及其上游表示学习模块理解 POI 隐含的地理信息，尤其是类别名称与真实地理位置之间的联合关系。围绕这一目标，本文先后考虑了显式注入与隐式注入两类方案。

显式注入方案的基本思想是直接优化 Generative POI Recommendation 模块，即在用户历史行为序列构造阶段，将地理信息以显式 token 的形式直接注入大语言模型的提示词中。具体而言，可以利用 Geohash 等地理编码方式将 POI 和用户的位置转换为统一的地理 token 表示，再与 SID 一同构成输入或输出目标。这样做的直观动机在于，希望模型在生成目标 POI SID 的同时，也能显式感知甚至生成对应的地理标识，从而强化空间信息利用。然而，实验过程中发现，该方案虽然从形式上看加强了地理信息暴露，但也显著增加了输出序列长度和生成复杂度。模型不仅要预测下一个 POI 的语义标识，还要同时处理额外的地理 token，导致训练与推理开销明显增加，推理速度甚至接近翻倍。同时，地理 token 与 SID 并列作为输出目标时，模型面临更大的建模压力，反而可能削弱对核心 SID 预测任务的学习效果，最终表现为准确率提升有限甚至下降。

与显式注入方案相比，本文进一步发现，将地理信息隐式注入 Semantic ID Construction 模块通常能够取得更好的效果。其核心思想是：不要求大语言模型显式生成地理标识，而是让地理信息在 POI 表示学习阶段直接参与连续表示构造，从而影响最终学得的 SID。这样一来，模型在输出阶段仍只需预测统一的 SID 序列，既避免了输出空间膨胀，又能够使地理信息真实地融入离散语义表示中。基于这一认识，本文最终选择通过 GeoPE 与文本编码器构造地理增强表示，并将其输入 RQ-VAE 学习地理增强 SID。

\section{显式地理信息注入方案及其局限性}
\label{sec:explicit_geo_attempt}

在方法探索初期，本文曾尝试通过显式方式将地理信息注入 Generative POI Recommendation 模块。该方案的基本出发点是：既然 GNPR-SID 最终通过大语言模型生成下一兴趣点的 SID，那么可以进一步将地理信息显式组织为 token，并直接加入输入提示词，甚至将其作为输出目标的一部分共同生成。为了统一表示 POI 与用户的位置，本文考虑过采用 Geohash 编码，将经纬度映射为离散地理串，再将其嵌入到用户访问历史和目标预测中。

在这一方案下，历史序列中的每个 POI 不再仅被表示为对应的 SID token，而是会进一步附加其地理编码。例如，在原始输入模板
\begin{quote}
\small
\texttt{user\_<uid> visited: <SID> at [time], ..., visited <SID> at [time].}
\end{quote}
的基础上，可以扩展为同时包含语义 token 和地理 token 的形式。对应地，模型输出也可以被设计为目标 POI 的“地理 token + SID token”联合结果。这样设计的优点在于，大语言模型可以在显式可见的 token 空间中直接学习 POI 的空间结构和区域相似性，从形式上更符合“地理信息参与生成”的直觉。

然而，实际实验表明，该方案存在明显局限。首先，地理 token 的显式加入会直接增加输入与输出序列长度。对于生成式模型而言，序列长度增加意味着自回归解码步骤增加，从而带来更高的训练与推理开销。尤其在 Next POI 推荐任务中，历史序列本身已经包含多个时间步访问记录，若每个 POI 都进一步附加地理编码，则输入会显著变长，而输出目标也会因联合生成地理 token 而变得更复杂。其次，地理 token 和 SID token 并列出现在输出端时，模型需要同时学习空间预测与语义预测两个子任务，这使训练目标更为分散。当地理信息不能有效帮助 SID 判别时，这种联合输出反而会稀释模型对核心预测目标的关注，导致准确率下降。

从实验观察来看，显式注入方案的一个重要问题在于：地理信息虽然被“看见”了，但未必真正转化为对核心 SID 预测的有效帮助。也就是说，模型可能能够学习一定程度的地理 token 模式，却无法将这种空间模式进一步转化为对下一个 POI 语义标识的判别能力。与此同时，地理 token 还增加了输出难度和推理复杂度，使整体收益不明显。因此，本文最终没有采用将地理信息直接并入输出目标的方式，而是转向更强调“通过表示学习隐式注入地理信息”的方案。

\section{基于 GeoPE 的隐式地理增强表示}
\label{sec:implicit_geope}

与显式注入方案相比，隐式注入方案更关注在连续表示层面建模地理信息，并通过这种连续表示影响离散 SID 的学习结果。本文借鉴了地理增强表示构造的相关思想，在 POI 表示学习阶段引入 GeoPE（Geographical Positional Encoding）机制，使经纬度信息能够以结构化方式隐式保存在潜在表示空间中。

设某一 POI 的文本属性为 $A_p$，在本文当前数据条件下，$A_p$ 主要对应于类别名称信息。由于数据集中缺少更丰富的文本描述，例如 POI 标题、地址或评论等，因此本文采用轻量级文本编码器对类别名称进行语义编码。在具体实现中，本文使用预训练模型 \texttt{all-MiniLM-L6-v2} 对类别文本进行编码，得到基础 POI 嵌入：
\begin{equation}
x_p = f_{text}(A_p),
\end{equation}
其中，$x_p \in \mathbb{R}^{D}$ 为文本编码器输出的连续表示。该模型输出维度约为 384 维，能够在较低计算成本下提供稳定的类别语义向量表示。之所以不采用更大规模的文本模型，是因为本文所使用的数据集中 POI 文本属性较为有限，仅类别名称本身不足以支撑复杂的大模型语义建模；相比之下，轻量文本编码器更符合数据规模与计算成本的平衡要求。

在获得基础语义表示 $x_p$ 后，本文进一步考虑如何将地理坐标信息注入该表示。GeoPE 的核心思想是，不直接将经纬度作为附加数值简单拼接，而是通过 POI 相对于若干参考点的方向关系，对嵌入向量进行旋转变换，从而使不同地理位置的 POI 在潜在表示空间中产生系统性差异。

具体而言，首先基于所有 POI 的地理分布，通过 K-Means 聚类得到 $\Omega$ 个参考点：
\begin{equation}
R=\{R_1,R_2,\dots,R_\Omega\}.
\end{equation}
每个参考点可视为一个区域密度中心，用于刻画整个城市空间中的若干典型区域。对于任意 POI $p$ 和参考点 $R_\omega$，计算二者之间的方位角
\begin{equation}
\theta_{p,\omega},
\end{equation}
该角度用于表示 POI 相对于参考点的方向关系。由于一个二维坐标点相对多个参考点的角度关系在几何上携带了充分的位置信息，因此这些角度可以被看作对真实地理位置的一种隐式编码。

随后，将基础嵌入 $x_p$ 划分为 $\Omega$ 个子段，每个子段对应一个参考点。对于第 $\omega$ 个子段中的相邻维度对，引入由方位角 $\theta_{p,\omega}$ 控制的二维旋转变换。设该子段中某一对相邻维度为 $(x_{2i}, x_{2i+1})$，则其旋转结果可表示为
\begin{equation}
\begin{bmatrix}
x'_{2i} \\
x'_{2i+1}
\end{bmatrix}
=
\begin{bmatrix}
\cos\theta_{p,\omega} & -\sin\theta_{p,\omega} \\
\sin\theta_{p,\omega} & \cos\theta_{p,\omega}
\end{bmatrix}
\begin{bmatrix}
x_{2i} \\
x_{2i+1}
\end{bmatrix}.
\end{equation}
对所有子段完成旋转后，再将其拼接得到最终的地理增强表示
\begin{equation}
x'_p = \mathrm{GeoPE}(x_p, l_p),
\end{equation}
其中 $l_p=(lat_p, lon_p)$ 表示 POI 的经纬度坐标。

通过上述方式，原本仅包含类别语义的基础嵌入 $x_p$ 被映射为同时包含语义属性与空间属性的增强表示 $x'_p$。与直接拼接经纬度相比，GeoPE 更强调通过变换方式将地理信息注入潜在空间，使空间差异成为连续表示内部结构的一部分。这种方式的优势在于：不同地理位置的 POI 即便属于同一类别，也会因相对参考点角度不同而得到不同的旋转结果，从而使模型能够在连续表示层面对同类异地 POI 进行更细致的区分。

\section{基于地理增强表示的 SID 学习}
\label{sec:geo_sid_learning}

在获得地理增强表示 $x'_p$ 后，本文将其作为 RQ-VAE 的输入，用于学习地理增强的离散语义标识 SID。与原始 GNPR-SID 中使用
\begin{equation}
p_e = \mathrm{concat}(c,r,t,c_u)
\end{equation}
作为输入不同，本文更加关注 POI 静态属性的地理--语义联合表达，因此采用由文本编码器与 GeoPE 共同构造的连续表示替代原始类别与区域文本输入。这样做的目标在于使后续量化过程能够直接作用于包含空间信息的潜在表示，从而使学得的 SID 更具地理判别能力。

设编码器输出的潜在表示为
\begin{equation}
z_e = Enc(x'_p),
\end{equation}
则 RQ-VAE 通过多级残差量化逐层逼近该表示，并在每一层选择与当前残差最接近的 codeword。若模型包含 $L$ 层码本，则每个 POI 最终对应一个由多级索引组成的离散表示：
\begin{equation}
SID(p)=[s_1,s_2,\dots,s_L].
\end{equation}
由于输入表示 $x'_p$ 中已经隐式融入了地理信息，因此该 SID 不再是单纯意义上的语义 ID，而是一种包含空间感知能力的离散语义标识。

本文仍然采用重建损失、量化损失与多样性损失共同优化 RQ-VAE，其总体损失形式可写为
\begin{equation}
L_{total}=L_{recon}+\mu L_{quant}+\lambda L_{div}.
\end{equation}
其中，$L_{recon}$ 用于保证量化后的表示仍能重构原始输入特征，$L_{quant}$ 用于稳定编码器输出与码本向量之间的映射关系，$L_{div}$ 用于鼓励码本利用更加均衡。通过该目标函数优化，模型能够逐步学习一组既保留语义结构又反映空间差异的离散标识。

需要强调的是，本文方法并不要求模型显式生成地理编码，而是让地理信息通过连续表示构造与量化学习过程隐式地体现在 SID 中。这种设计的核心优势在于：输出阶段仍然保持统一的 SID 预测形式，不会因引入额外的地理 token 而显著增加生成复杂度；与此同时，SID 的学习结果又能够受益于地理增强表示，从而在不改变生成任务形式的前提下提升 POI 表示能力。

\section{地理增强 SID 在生成式 POI 推荐中的应用}
\label{sec:geo_sid_application}

在完成地理增强 SID 学习后，本文沿用 GNPR-SID 的生成式推荐框架，将每个 POI 对应的 SID 组织为结构化 token 序列，并用于构造用户历史访问序列的输入提示。设某一 POI 的 SID 为
\begin{equation}
SID(p)=[s_1,s_2,\dots,s_L],
\end{equation}
则可进一步映射为
\begin{equation}
T(p)=\langle a_{s_1}\rangle \langle b_{s_2}\rangle \cdots \langle l_{s_L}\rangle.
\end{equation}

对于用户 $u$ 的历史访问序列
\begin{equation}
H_u=\{(p_1,t_1),(p_2,t_2),\dots,(p_n,t_n)\},
\end{equation}
本文将其组织为如下形式的输入提示：
\begin{quote}
\small
\texttt{user\_<uid> visited: <SID> at [time], ..., visited <SID> at [time]. When [time] user\_<uid> is likely to visit:}
\end{quote}
与原始 GNPR-SID 的不同之处在于，这里的 \texttt{<SID>} 不再仅来自类别、区域与协同特征学习得到的语义 ID，而是来源于地理增强表示 $x'_p$ 经 RQ-VAE 学习得到的地理增强 SID。因此，尽管生成式模型在输出端仍只预测统一的 SID 序列，但其预测目标本身已经隐含了更多空间信息。

在训练阶段，本文采用 supervised fine-tuning 的方式对大语言模型进行微调，使其根据用户历史访问记录和目标时间上下文生成下一 POI 的 SID。在预测阶段，模型输出的 SID 再映射回对应 POI，从而完成 Next POI 推荐任务。

\section{本章小结}
\label{sec:method_summary}

本章围绕本文提出的基于地理信息增强的 GNPR-SID 方法进行了详细介绍。首先，对 GNPR-SID 原始框架进行了梳理，说明其由 Semantic ID Construction 和 Generative POI Recommendation 两部分组成。随后，分析了原始方法在地理信息建模方面的不足，并对显式注入地理信息的方案及其局限进行了讨论。基于实验现象和方法分析，本文最终选择在 Semantic ID Construction 阶段隐式注入地理信息，通过轻量文本编码器和 GeoPE 构造地理增强表示，并将其输入 RQ-VAE 学习地理增强 SID。最后，说明了该 SID 如何在后续生成式 POI 推荐模型中使用。

与直接生成地理标识相比，本文方法更强调让地理信息影响 SID 的学习结果，而不是将其作为额外输出目标。这一设计既保持了 GNPR-SID 原有生成式推荐框架的统一性，又提升了离散语义表示的空间感知能力，为后续实验验证提供了方法基础。